% GENERATED FILE. Edit canonical lessons, then run npm run generate:books.
% canonical-source-hash: fnv1a64:ba7330e58286dda7
% canonical-lessons: AR-C396-mustalah, AR-C396-tarif, AR-C396-fihris, AR-C396-utruha, AR-C396-manhaj

\chapter{Terms and theses}
\label{ch:ar-terms-and-theses}
\hlchaptermodality{396}

\begin{chapteropening}
\textbf{By the end of this chapter:} \emph{I can talk about a term, a definition, an index, a thesis and a curriculum.}

Complete the last lesson of chapter 396: I can talk about a term, a definition, an index, a thesis and a curriculum.
\end{chapteropening}

\section[\texorpdfstring{\ar{مصطلح} (muṣṭalaḥ)}{مصطلح (muṣṭalaḥ)}]{\ar{مصطلح} (muṣṭalaḥ) — a term}

\label{lesson:AR-C396-mustalah}



\begin{quote}
Before the new one: say the Arabic for to read up on, then the Arabic for to reflect.
\end{quote}

\subsection*{You'll want to know: \ar{مصطلح}}
\textbf{\ar{مصطلح}} — \emph{muṣṭalaḥ} — \textquotedblleft{}a term\textquotedblright{}.

\begin{quote}
\ar{هذا} \ar{مصطلح} \ar{علمي} \ar{جديد}.
\end{quote}

\begin{quote}
\emph{hādhā muṣṭalaḥ ʿilmī jadīd.}
\end{quote}

\begin{quote}
This is a new scientific term.
\end{quote}

\begin{cousinweb}
\textbf{\ar{مصطلح}} (\emph{muṣṭalaḥ}) is built on the native Arabic root \emph{ṣ-l-ḥ}, \textquotedblleft{}to be good, right\textquotedblright{} (Lane). \emph{Muṣṭalaḥ} is the Form VIII passive participle, what people have agreed on: a term.
\end{cousinweb}

\begin{grammarlens}[title={What you've built}]
One more word for giving reasons and opinions.
\end{grammarlens}

\subsection*{Your turn}
\begin{itemize}
\raggedright
  \item \emph{Say it:} \emph{muṣṭalaḥ}
  \item \emph{Say it:} \emph{muṣṭalaḥ}, once more
  \item \emph{Say it:} \emph{hādhā muṣṭalaḥ ʿilmī jadīd.} — This is a new scientific term.
  \item \emph{Recall:} say the Arabic for to read up on, then the Arabic for to reflect, then say \emph{muṣṭalaḥ} again
\end{itemize}

\begin{tcolorbox}[breakable,colback=teal!4,colframe=teal!35!black,title={Before you move on}]
What does \ar{مصطلح} mean? (\textquotedblleft{}A term\textquotedblright{}.) Say it once more.
\end{tcolorbox}
\section[\texorpdfstring{\ar{تعريف} (taʿrīf)}{تعريف (taʿrīf)}]{\ar{تعريف} (taʿrīf) — a definition}

\label{lesson:AR-C396-tarif}



\begin{quote}
Before the new one: say the Arabic for to reflect, then the Arabic for a term.
\end{quote}

\subsection*{You'll want to know: \ar{تعريف}}
\textbf{\ar{تعريف}} — \emph{taʿrīf} — \textquotedblleft{}a definition\textquotedblright{}.

\begin{quote}
\ar{كتب} \ar{المعلم} \ar{تعريف} \ar{الكلمة}.
\end{quote}

\begin{quote}
\emph{kataba al-muʿallim taʿrīf al-kalima.}
\end{quote}

\begin{quote}
The teacher wrote down the definition of the word.
\end{quote}

\begin{cousinweb}
\textbf{\ar{تعريف}} (\emph{taʿrīf}) is built on the native Arabic root \emph{ʿ-r-f}, \textquotedblleft{}to know\textquotedblright{} (Lane). \emph{Taʿrīf} is the Form II verbal noun, making known: a definition.
\end{cousinweb}

\begin{grammarlens}[title={What you've built}]
One more word for giving reasons and opinions.
\end{grammarlens}

\subsection*{Your turn}
\begin{itemize}
\raggedright
  \item \emph{Say it:} \emph{taʿrīf}
  \item \emph{Say it:} \emph{taʿrīf}, once more
  \item \emph{Say it:} \emph{kataba al-muʿallim taʿrīf al-kalima.} — The teacher wrote down the definition of the word.
  \item \emph{Recall:} say the Arabic for to reflect, then the Arabic for a term, then say \emph{taʿrīf} again
\end{itemize}

\begin{tcolorbox}[breakable,colback=teal!4,colframe=teal!35!black,title={Before you move on}]
What does \ar{تعريف} mean? (\textquotedblleft{}A definition\textquotedblright{}.) Say it once more.
\end{tcolorbox}
\section[\texorpdfstring{\ar{فهرس} (fihris)}{فهرس (fihris)}]{\ar{فهرس} (fihris) — an index}

\label{lesson:AR-C396-fihris}



\begin{quote}
Before the new one: say the Arabic for a term, then the Arabic for a definition.
\end{quote}

\subsection*{You'll want to know: \ar{فهرس}}
\textbf{\ar{فهرس}} — \emph{fihris} — \textquotedblleft{}an index\textquotedblright{}.

\begin{quote}
\ar{وجدت} \ar{الكلمة} \ar{في} \ar{فهرس} \ar{الكتاب}.
\end{quote}

\begin{quote}
\emph{wajadtu al-kalima fī fihris al-kitāb.}
\end{quote}

\begin{quote}
I found the word in the book's index.
\end{quote}

Plural: \ar{فهارس}, \emph{fahāris}.

\begin{cousinweb}
\textbf{\ar{فهرس}} (\emph{fihris}) was borrowed from Persian \emph{fihrist}, \textquotedblleft{}a list\textquotedblright{} (EtymDB); Lane calls it an arabicized word for a catalogue of books (Lane). Today it is an index or table of contents.
\end{cousinweb}

\begin{grammarlens}[title={What you've built}]
One more word for giving reasons and opinions.
\end{grammarlens}

\subsection*{Your turn}
\begin{itemize}
\raggedright
  \item \emph{Say it:} \emph{fihris}
  \item \emph{Say it:} \emph{fihris}, once more
  \item \emph{Say it:} \emph{wajadtu al-kalima fī fihris al-kitāb.} — I found the word in the book's index.
  \item \emph{Recall:} say the Arabic for a term, then the Arabic for a definition, then say \emph{fihris} again
\end{itemize}

\begin{tcolorbox}[breakable,colback=teal!4,colframe=teal!35!black,title={Before you move on}]
What does \ar{فهرس} mean? (\textquotedblleft{}An index\textquotedblright{}.) Say it once more.
\end{tcolorbox}
\section[\texorpdfstring{\ar{أطروحة} (uṭrūḥa)}{أطروحة (uṭrūḥa)}]{\ar{أطروحة} (uṭrūḥa) — a thesis}

\label{lesson:AR-C396-utruha}



\begin{quote}
Before the new one: say the Arabic for a definition, then the Arabic for an index.
\end{quote}

\subsection*{You'll want to know: \ar{أطروحة}}
\textbf{\ar{أطروحة}} — \emph{uṭrūḥa} — \textquotedblleft{}a thesis\textquotedblright{}.

\begin{quote}
\ar{كتب} \ar{الباحث} \ar{أطروحة} \ar{عن} \ar{الهجرة}.
\end{quote}

\begin{quote}
\emph{kataba al-bāḥith uṭrūḥa ʿan al-hijra.}
\end{quote}

\begin{quote}
The researcher wrote a thesis on migration.
\end{quote}

\begin{cousinweb}
\textbf{\ar{أطروحة}} (\emph{uṭrūḥa}) is built on the native Arabic root \emph{ṭ-r-ḥ}, \textquotedblleft{}to throw, to cast\textquotedblright{} (Lane). An \emph{uṭrūḥa} is a thing put forward, thrown out for discussion: a thesis.
\end{cousinweb}

\begin{grammarlens}[title={What you've built}]
One more word for giving reasons and opinions.
\end{grammarlens}

\subsection*{Your turn}
\begin{itemize}
\raggedright
  \item \emph{Say it:} \emph{uṭrūḥa}
  \item \emph{Say it:} \emph{uṭrūḥa}, once more
  \item \emph{Say it:} \emph{kataba al-bāḥith uṭrūḥa ʿan al-hijra.} — The researcher wrote a thesis on migration.
  \item \emph{Recall:} say the Arabic for a definition, then the Arabic for an index, then say \emph{uṭrūḥa} again
\end{itemize}

\begin{tcolorbox}[breakable,colback=teal!4,colframe=teal!35!black,title={Before you move on}]
What does \ar{أطروحة} mean? (\textquotedblleft{}A thesis\textquotedblright{}.) Say it once more.
\end{tcolorbox}
\section[\texorpdfstring{\ar{منهج} (manhaj)}{منهج (manhaj)}]{\ar{منهج} (manhaj) — a curriculum}

\label{lesson:AR-C396-manhaj}



\begin{quote}
Before the new one: say the Arabic for an index, then the Arabic for a thesis.
\end{quote}

\subsection*{You'll want to know: \ar{منهج}}
\textbf{\ar{منهج}} — \emph{manhaj} — \textquotedblleft{}a curriculum\textquotedblright{}.

\begin{quote}
\ar{غيرت} \ar{المدرسة} \ar{منهج} \ar{العلوم}.
\end{quote}

\begin{quote}
\emph{ghayyarat al-madrasa manhaj al-ʿulūm.}
\end{quote}

\begin{quote}
The school changed the science curriculum.
\end{quote}

Plural: \ar{مناهج}, \emph{manāhij}.

\begin{cousinweb}
\textbf{\ar{منهج}} (\emph{manhaj}) is built on the native Arabic root \emph{n-h-j}, \textquotedblleft{}to be plain, open (of a road)\textquotedblright{} (Lane). A \emph{manhaj} is a plain road: a method, and in schools a curriculum.
\end{cousinweb}

\begin{grammarlens}[title={What you've built}]
One more word for giving reasons and opinions.
\end{grammarlens}

\subsection*{Your turn}
\begin{itemize}
\raggedright
  \item \emph{Say it:} \emph{manhaj}
  \item \emph{Say it:} \emph{manhaj}, once more
  \item \emph{Say it:} \emph{ghayyarat al-madrasa manhaj al-ʿulūm.} — The school changed the science curriculum.
  \item \emph{Recall:} say the Arabic for an index, then the Arabic for a thesis, then say \emph{manhaj} again
\end{itemize}

\begin{tcolorbox}[breakable,colback=teal!4,colframe=teal!35!black,title={Before you move on}]
What does \ar{منهج} mean? (\textquotedblleft{}A curriculum\textquotedblright{}.) Say it once more.
\end{tcolorbox}
